\section{Conifold geometry and actual metric comparison}\label{sec:geometry}

\subsection{The link and its integral normalization}
An ordinary double point is analytically the cone $\sum_{j=1}^4z_j^2=0$.
Its Ricci-flat metric cone is $g_C=dr^2+r^2h$, where the homogeneous
metric on $T^{1,1}\simeq S^2\times S^3$ is
\begin{equation}\label{eq:link-metric}
 h=\frac19\eta_c^2+
 \frac16\sum_{j=1}^2(d\theta_j^2+\sin^2\theta_j\,d\phi_j^2),\qquad
 \eta_c=d\psi+\cos\theta_1d\phi_1+\cos\theta_2d\phi_2.
\end{equation}
Here $\psi$ has period $4\pi$ and $\Ric_h=4h$.
Put $\sigma_j=\sin\theta_jd\theta_j\wedge d\phi_j$ and
\begin{equation}\label{eq:link-form}
 \vartheta=\frac{\sigma_1-\sigma_2}{8\pi}.
\end{equation}
The circle bundle over $\PP^1\times\PP^1$ has Chern class $(1,1)$ up
to a common sign. A primitive two-cycle projects with degrees $(1,-1)$;
therefore $\vartheta$ has period one. It is harmonic: it is closed,
its horizontal part is anti-self-dual, and wedging that part with
$d\eta_c=-(\sigma_1+\sigma_2)$ gives zero. In the orthonormal frame,
\begin{equation}\label{eq:link-norm}
 |\vartheta|_h^2=\frac9{8\pi^2},\qquad
 \vol(T^{1,1},h)=\frac{16\pi^3}{27},\qquad
 c:=\|\vartheta\|_{L^2(h)}^2=\frac{2\pi}{3}.
\end{equation}
The selected small-resolution ruling fixes the sign of its primitive
period. On a cone, $\vartheta$ is a primitive real $(1,1)$-form, of
pointwise size $O(r^{-2})$.

The smoothing cap is $T^*S^3$ with the Stenzel metric, asymptotic with
positive differentiated rate to this cone \cite{stenzel,candelas-ossa}.
The small-resolution cap is the total space of
$\OO_{\PP^1}(-1)^{\oplus2}$ with its Ricci-flat resolved-conifold metric.
The two caps have the same link but retain different primitive cohomology:
restriction to the link is an isomorphism in degree three for the smoothing
cap and in degree two for the resolution cap.

\subsection{The central metric and the smoothing scale}
Let $X_0^\reg$ carry the singular Ricci-flat K\"ahler form $\omega$ in
class $H$. The bounded-potential construction \cite[\S7.2, Theorem~7.5 and Corollary~7.6]{egz} gives a closed positive pullback current on $Y$. Its
Bedford--Taylor products have their cohomological masses and give no mass
to the exceptional analytic set. The conical theorem of Hein--Sun
\cite[Theorem~1.4 and Corollary~1.5]{hein-sun} supplies holomorphic node
charts such that, for some $\mu>0$,
\begin{equation}\label{eq:central-rate}
 |\nabla_C^j(g_0-g_C)|_{g_C}\le C_jr^{\mu-j}
 \quad (j\ge0).
\end{equation}
Use the minimum rate over the finitely many ends. Benabida's Theorem~A
provides the corresponding polyhomogeneous expansion \cite{benabida}.
The cited metric theorems use the class $2\pi c_1(\mathcal L)$ for
the polarizing line bundle $\mathcal L$. We rescale their metrics by
$(2\pi)^{-1}$ to obtain our class $H=c_1(\mathcal L)$, with the corresponding fixed rescaling
of cone radii and cap normalizations. Written as $dr^2+r^2h$, the
normalized link metric and \eqref{eq:link-norm} are unchanged.

Condition~\eqref{eq:simple-smoothing} makes the total-space projection a
holomorphic Morse function near each node. The local family therefore
has the weighted ordinary-node model, with weights $3/2$ on the coordinates
and $3$ on the algebraic smoothing parameter. This model is compatible
with the central conical chart. Indeed, the transition between two
ordinary-quadric charts lifts to an ambient holomorphic germ $\Phi$ with
$Q\circ\Phi=uQ$, where $Q=\sum z_j^2$ and $u$ is a holomorphic unit.
Its linear part $A$ satisfies $Q(Az)=u(0)Q(z)$ and extends to the
weighted smoothing after a fixed rescaling of the base. Since all
coordinate weights are $3/2$, the terms $\Phi-A=O(|z|^2)$ have relative
conic order $O(r^{3/2})$, with differentiated bounds. After this linear
normalization, \eqref{eq:central-rate} therefore holds in the family
chart, with $\mu$ decreased if necessary. Along the fixed positive ray
we use the reference metric scale
\begin{equation}\label{eq:metric-scale}
 \epsilon=t^{1/3}.
\end{equation}
On the smoothing branch, metrics, K\"ahler forms and eigenvalues indexed
by $\epsilon$ are evaluated at the algebraic parameter $t=\epsilon^3$.
The individual cap scales have fixed positive leading ratios to $\epsilon$.
Multiplying a local normalization by a fixed constant changes a logarithmic
length by a bounded term.

\begin{lemma}[Uniform corner extraction]\label{lem:corner-extraction}
Let $f(\rho_1,\rho_2,y)$ be a polyhomogeneous tensor coefficient on a
compact product corner. Suppose its index sets at both faces have only
log-free order-zero terms and terms of strictly positive real order, and
let $f_0(\rho_1,y)$ be its trace at $\rho_2=0$. There is $a_0>0$ such
that, for every fixed number of b-derivatives,
\begin{equation}\label{eq:corner-extraction}
 |D_b^j(f-f_0)|\le C_j\rho_2^{a_0}
\end{equation}
uniformly up to $\rho_1=0$.
\end{lemma}
\begin{proof}
Discreteness of the index set gives a positive gap above zero at the
$\rho_2$ face. Choose $a_0$ below that gap and a positive remainder order.
Expansion coefficients and their fixed b-derivatives are bounded at the
other face, since its index set has no negative orders or order-zero
logarithms. Subtracting the trace removes the order-zero term; positive
powers absorb the remaining finite logarithmic powers. The conormal
remainder has the same bound. A finite corner cover makes the constants
uniform. Pointwise convergence at fixed $\rho_1>0$ alone would not imply
this conclusion.
\end{proof}

\begin{proposition}[Actual metric comparison]\label{prop:metric-comparison}
There are collar-compatible exterior identifications $F_\epsilon$ and
constants $a_0>0$, $r_0>0$ such that, for fixed sufficiently large $A$,
\begin{equation}\label{eq:actual-comparison}
 |\nabla_C^j(F_\epsilon^*g_\epsilon-g_0)|_{g_C}
 \le C_j(\epsilon/r)^{a_0}r^{-j},
 \quad A\epsilon\le r\le r_0,\quad j=0,1.
\end{equation}
The same comparison holds for the K\"ahler forms. On fixed exterior
compact sets the $C^1$ differences are $O(\epsilon^{a_0})$.
On fixed compact subsets of each rescaled cap, the metrics converge in
all derivatives to the normalized Stenzel metric.
\end{proposition}

\begin{proof}
We use the construction of \cite[Theorem~C]{benabida}, with the following
finite-node extension and tensor deduction. The cited work is published;
our numbered references to its proof refer to arXiv version~3, as specified
in the bibliography. Remark~5.2 there describes the extension, but its
Lemma~5.4 is formulated with one cap face. We give the global estimate
needed to replace that hypothesis.

\emph{The normalized scalar inverse.}
Let $\rho_\epsilon$ be a positive radius, equal to the central radius on
the necks and comparable to $\epsilon$ on the cores. For $-4<\beta<0$ put
\[
 \|u\|_{H_\beta^q}^2=
 \sum_{j=0}^q\int |\rho_\epsilon^{j-\beta}\nabla^ju|^2
                         \rho_\epsilon^{-6}\,dV.
\]
The CS space in \cite[Definitions~11.1 and~11.6]{pacini} is the connected
central regular locus; the AC space is the disjoint union of the normalized
Stenzel caps. Disconnected AC spaces are allowed. All ends are filled at
scale $\epsilon$, with fixed ratios absorbed into the caps. Choose a smooth
density $\zeta$ of integral one supported in a fixed central compact set.
Theorem~12.3(1) of \cite{pacini} gives the estimate
\begin{equation}\label{eq:uniform-scalar-inverse}
 \|u\|_{H_\beta^{q+2}}\le C_q\|\Delta_\epsilon u\|_{H_{\beta-2}^q},
 \qquad u\in E_\epsilon:=\left\{\int\zeta u=0\right\},
\end{equation}
uniformly in $\epsilon$, for its parametric connected-sum metric. The
complement is precisely the pullback of the fixed central complement in
that proof: the cutoff used there is one on $\operatorname{supp}\zeta$.
The restricted Laplacian has the same image as the full Laplacian. On the
connected closed smoothing this image is
$\{f:\int f\,dV_\epsilon=0\}$, and the kernel is the single space of
constants. Thus the range condition is global; there is no separate
normalization or mean-zero condition for each cap.

The initial K\"ahler tensor in \cite[Proposition~5.9]{benabida} has the
prescribed central and cap traces with positive orders at the opposite
faces. Compare it with Pacini's metric, interpolated at
$r=\epsilon^\tau$, $0<\tau<1$. On the central and cap sides, its initial
face expansions give errors bounded by positive powers of
$\epsilon^{1-\tau}$ and $\epsilon^\tau$. On the intervening annuli,
\eqref{eq:central-rate} and the AC rate give the same conclusion in every
fixed scaled differentiability norm. Therefore
\[
 \|\Delta_{\rm initial}-\Delta_{\rm Pacini}\|
       _{H_\beta^{q+2}\to H_{\beta-2}^q}=o(1).
\]
Absorption in \eqref{eq:uniform-scalar-inverse} proves the estimate for the
initial family and its subsequent boundary-vanishing formal corrections.
This comparison uses the initial compactification, not any estimate for
the final Ricci-flat metric. Theorem~11.7 of \cite{pacini} supplies the
uniform scaled Sobolev and multiplication estimates used in the nonlinear
argument.

\emph{Finite-node formal corrections.}
The surgery space has one connected central face and finitely many pairwise
disjoint cap faces. The central coefficient problems of
\cite[\S\S5--6]{benabida} are scalar equations on that one connected
conifold, with one integral compatibility condition. The finite-end
positive-weight inverse includes the low harmonic terms at every end
\cite[Proposition~3.4]{benabida}. These terms cannot be discarded when
differentiating the potentials. For the node, the first scalar link
eigenvalues are $0,33/4,12$, with growth rates $0,3/2,2$; below growth two
the terms are constants and the real and imaginary parts of the four
holomorphic linear coordinates \cite[Theorem~2.14(1), Example~2.16]{hein-sun}.
Extend the corresponding central holomorphic germs through the smoothing
charts and multiply by disjoint cutoffs equal to one near the nodes.
Their fibrewise $\partial\bar\partial$ vanishes near the corners, and
all cutoff derivatives have fixed exterior support. Consequently the low
terms introduce no negative corner tensor order.

The AC coefficient equations are a finite direct sum of the inverses in
\cite[Proposition~3.6]{benabida}. Their decaying weight range has no
componentwise integral obstruction. The global volume normalization gives
the compatibility for the central equation. More explicitly, each cap's
volume starts at order $\epsilon^6$, so a central residual coefficient
of order less than six has zero integral by the push-forward argument of
\cite[Proposition~5.5]{benabida}. At the next stage one improves the
central residual to almost order six and the cap residuals to almost order
four, factors out $\epsilon^{4-\eta}$ for a sufficiently small $\eta>0$,
and repeats the same procedure. Multiplication by this scalar preserves
the global zero-integral identity. Cap contributions add in that identity;
they are not required to vanish individually. Finitely many disjoint faces
permit common choices of weights and the same Borel summation.

At a product corner write $\epsilon=\rho_1\rho_2$. The tensor orders in
\cite[proof of Proposition~6.1]{benabida} are the following. A central
potential correction $\epsilon^\delta w$, after removing the
pluriharmonic terms just treated, has $w$ of weight $a>2-\delta$.
Choose $\max(\delta/2,2-a)<\delta'<\delta$. Its twice-differentiated
conic tensor has order
\[
 \rho_2^{\delta'}\rho_1^{a-2+\delta'}.
\]
Both exponents are positive, logarithms are absorbed by the loss in
$\delta'$, and $2\delta'>\delta$ improves the nonlinear residual.
A cap correction has tensor order
$\rho_1^{\delta'}\rho_2^{a+2+\delta'}$, with positive exponents;
the cap weight choice preserves the previously improved central order.
These are tensor estimates, not merely positive orders for potentials.

The inverse \eqref{eq:uniform-scalar-inverse} and the scaled nonlinear
bounds give the final correction as in \cite[\S7]{benabida}. With the
same normalization for every order, uniqueness and Sobolev embedding give
\[
 |\nabla^j\phi_\epsilon|\le C_{N,j}\epsilon^N
                           \rho_\epsilon^{\beta-j}
 \quad\text{for every fixed }N,j.
\]
Since $\epsilon=\rho_1\rho_2$ and $\rho_\epsilon$ is comparable to
$\rho_1$, its metric correction is flat at both faces after any fixed
number of conic derivatives. Joint flatness follows by differentiating
the normalized parameter equation: the coefficient derivatives cost only
fixed powers, and the same uniform linearized inverse gives the displayed
arbitrary-order decay inductively. Thus fibrewise estimates are not being
used alone to infer transverse regularity.

The final tensors consequently have nonnegative polyhomogeneous index
families in the vertical conic bundles: the order-zero term at each face
is log-free, and every other exponent has positive real part. The face
values asserted in Theorem~C are unchanged. The central value is the full
$g_0$; on a cap, $\epsilon^{-2}g_\epsilon$ restricts to its normalized
Stenzel metric. In the conic bundle the latter is
$\rho_2^2g_{\rm St}$ in the cap b-frame. Since
$\rho_2=(1+\widetilde r^2)^{-1/2}$, it is bounded up to the corner
and has the cone trace there.

\emph{Extraction in the conic frame.}
The index families just obtained satisfy Lemma~\ref{lem:corner-extraction}.
Map $(y,r,0)$ to $(y,r,\epsilon/r)$, taking $r=\rho_1$. The vertical
b-frame is $\rho_1\partial_{\rho_1}-\rho_2\partial_{\rho_2},\partial_y$.
After division by $\rho_1$ its pullback is the cone frame
$\partial_r,r^{-1}\partial_y$. A cone covariant derivative costs $r^{-1}$
and bounded b-differentiation. Equation~\eqref{eq:corner-extraction} gives
\eqref{eq:actual-comparison}. Extend the prescribed collar lifts by a
partition of unity on the smooth fixed exterior; their flow gives compatible
identifications. There is no holomorphicity requirement. The exterior and
rescaled-cap estimates follow from the respective face restrictions. The
same argument applies to the K\"ahler tensors.

Finally, the initial form in \cite[Proposition~5.9]{benabida} is obtained
from curvature on the actual polarizing line bundle. The Stenzel form has
a global potential, and the gluing and all subsequent corrections are
fibrewise $\partial\bar\partial$-exact. The construction therefore stays
in $2\pi c_1(\mathcal L_t)$. Ricci-flat uniqueness in this class, followed
by the fixed rescaling already specified, identifies it with $g_t$ on the
prescribed complex family.
\end{proof}

The trace in \eqref{eq:corner-extraction} is the full central tensor.
Comparing both metrics separately with the cone would leave an
$O(r^\mu)$ term independent of $\epsilon$; that weaker estimate is not
used for the differentiated remainder argument.
